% GENERATED FILE. Edit canonical lessons, then run npm run generate:books.
% canonical-source-hash: fnv1a64:59f0e8f67e3c8007
% canonical-lessons: PA-C231-tel, PA-C231-gundna, PA-C231-gutt, PA-C231-paranda, PA-C231-lat, PA-R231-review-225-231

\chapter{Combing and Plaiting}
\label{ch:pa-a2-combing-and-plaiting}
\hlchaptermodality{231}

\begin{chapteropening}
\textbf{By the end of this chapter:} \emph{I can talk about doing my hair.}

Complete the last lesson of chapter 231: I can talk about doing my hair.
\end{chapteropening}

\section[\texorpdfstring{\pa{ਤੇਲ} (tel)}{ਤੇਲ (tel)}]{\pa{ਤੇਲ} (tel) — oil}

\label{lesson:PA-C231-tel}



\begin{quote}
Before the new one: say the Punjabi for loose trousers, then the Punjabi for a hanging earring.
\end{quote}

\subsection*{You'll want to know: \pa{ਤੇਲ}}
\textbf{\pa{ਤੇਲ}} — \emph{tel} — \textquotedblleft{}oil\textquotedblright{}.

\begin{quote}
\pa{ਮੈਂ} \pa{ਵਾਲਾਂ} \pa{ਉੱਤੇ} \pa{ਤੇਲ} \pa{ਮਲਦੀ} \pa{ਹਾਂ।}
\end{quote}

\begin{quote}
\emph{maiṁ vālāṁ utte tel maldī hāṁ.}
\end{quote}

\begin{quote}
I rub oil into my hair.
\end{quote}

A masculine noun.

\begin{cousinweb}
\textbf{\pa{ਤੇਲ}} (\emph{tel}) is inherited from Sanskrit \emph{taila}, \textquotedblleft{}sesame oil, oil\textquotedblright{}; CDIAL lists Punjabi \emph{tel}, \textquotedblleft{}oil\textquotedblright{} (CDIAL 5958).
\end{cousinweb}

\begin{grammarlens}[title={What you've built}]
One more word for saying what you do.
\end{grammarlens}

\subsection*{Your turn}
\begin{itemize}
\raggedright
  \item \emph{Say it:} \emph{tel}
  \item \emph{Say it:} \emph{tel}, once more
  \item \emph{Say it:} \emph{maiṁ vālāṁ utte tel maldī hāṁ.} — I rub oil into my hair.
  \item \emph{Recall:} say the Punjabi for loose trousers, then the Punjabi for a hanging earring, then say \emph{tel} again
\end{itemize}

\begin{tcolorbox}[breakable,colback=teal!4,colframe=teal!35!black,title={Before you move on}]
What does \pa{ਤੇਲ} mean? (\textquotedblleft{}Oil\textquotedblright{}.) Say it once more.
\end{tcolorbox}
\section[\texorpdfstring{\pa{ਗੁੰਦਣਾ} (gundṇā)}{ਗੁੰਦਣਾ (gundṇā)}]{\pa{ਗੁੰਦਣਾ} (gundṇā) — to braid (hair)}

\label{lesson:PA-C231-gundna}



\begin{quote}
Before the new one: say the Punjabi for a hanging earring, then the Punjabi for oil.
\end{quote}

\subsection*{You'll want to know: \pa{ਗੁੰਦਣਾ}}
\textbf{\pa{ਗੁੰਦਣਾ}} — \emph{gundṇā} — \textquotedblleft{}to braid (hair)\textquotedblright{}.

\begin{quote}
\pa{ਮਾਤਾ} \pa{ਸਵੇਰੇ} \pa{ਮੇਰੇ} \pa{ਵਾਲ} \pa{ਗੁੰਦਦੀ} \pa{ਹੈ।}
\end{quote}

\begin{quote}
\emph{mātā savere mere vāl gunddī hai.}
\end{quote}

\begin{quote}
My mother braids my hair in the morning.
\end{quote}

Past (I …): \pa{ਮੈਂ} \pa{ਵਾਲ} \pa{ਗੁੰਦੇ}, \emph{maiṁ vāl gunde}.

\begin{cousinweb}
\textbf{\pa{ਗੁੰਦਣਾ}} (\emph{gundṇā}) is inherited from Sanskrit \emph{guphati}, \textquotedblleft{}strings together\textquotedblright{}; CDIAL lists Punjabi \emph{gundṇā}, \textquotedblleft{}to braid hair\textquotedblright{} (CDIAL 4205).
\end{cousinweb}

\begin{grammarlens}[title={What you've built}]
One more word for saying what you do.
\end{grammarlens}

\subsection*{Your turn}
\begin{itemize}
\raggedright
  \item \emph{Say it:} \emph{gundṇā}
  \item \emph{Say it:} \emph{gundṇā}, once more
  \item \emph{Say it:} \emph{mātā savere mere vāl gunddī hai.} — My mother braids my hair in the morning.
  \item \emph{Recall:} say the Punjabi for a hanging earring, then the Punjabi for oil, then say \emph{gundṇā} again
  \item \emph{Write it:} copy \textbf{\pa{ਗੁੰਦਣਾ}} once, checking the whole word
\end{itemize}

\begin{tcolorbox}[breakable,colback=teal!4,colframe=teal!35!black,title={Before you move on}]
What does \pa{ਗੁੰਦਣਾ} mean? (\textquotedblleft{}To braid (hair)\textquotedblright{}.) Say it once more.
\end{tcolorbox}
\section[\texorpdfstring{\pa{ਗੁੱਤ} (gutt)}{ਗੁੱਤ (gutt)}]{\pa{ਗੁੱਤ} (gutt) — a plait of hair}

\label{lesson:PA-C231-gutt}



\begin{quote}
Before the new one: say the Punjabi for oil, then the Punjabi for to braid.
\end{quote}

\subsection*{You'll want to know: \pa{ਗੁੱਤ}}
\textbf{\pa{ਗੁੱਤ}} — \emph{gutt} — \textquotedblleft{}a plait of hair\textquotedblright{}.

\begin{quote}
\pa{ਕੁੜੀ} \pa{ਦੀ} \pa{ਗੁੱਤ} \pa{ਲੰਮੀ} \pa{ਹੈ।}
\end{quote}

\begin{quote}
\emph{kuṛī dī gutt lammī hai.}
\end{quote}

\begin{quote}
The girl's plait is long.
\end{quote}

A feminine noun.

\begin{cousinweb}
\textbf{\pa{ਗੁੱਤ}} (\emph{gutt}) comes from the same root as \pa{ਗੁੰਦਣਾ}, Sanskrit \emph{guphati}, \textquotedblleft{}strings together\textquotedblright{}: CDIAL lists Punjabi \emph{gundṇā}, \textquotedblleft{}to braid hair\textquotedblright{}, under it, but has no Punjabi row for the noun; it lists only Lahnda \emph{gutt}, \textquotedblleft{}plait of hair hanging down woman's back\textquotedblright{} (CDIAL 4205).
\end{cousinweb}

\begin{grammarlens}[title={What you've built}]
One more word for saying what you do.
\end{grammarlens}

\subsection*{Your turn}
\begin{itemize}
\raggedright
  \item \emph{Say it:} \emph{gutt}
  \item \emph{Say it:} \emph{gutt}, once more
  \item \emph{Say it:} \emph{kuṛī dī gutt lammī hai.} — The girl's plait is long.
  \item \emph{Recall:} say the Punjabi for oil, then the Punjabi for to braid, then say \emph{gutt} again
  \item \emph{Write it:} copy \textbf{\pa{ਗੁੱਤ}} once, checking the whole word
\end{itemize}

\begin{tcolorbox}[breakable,colback=teal!4,colframe=teal!35!black,title={Before you move on}]
What does \pa{ਗੁੱਤ} mean? (\textquotedblleft{}A plait of hair\textquotedblright{}.) Say it once more.
\end{tcolorbox}
\section[\texorpdfstring{\pa{ਪਰਾਂਦਾ} (parāndā)}{ਪਰਾਂਦਾ (parāndā)}]{\pa{ਪਰਾਂਦਾ} (parāndā) — a hair tassel; coloured yarn tied into a plait}

\label{lesson:PA-C231-paranda}



\begin{quote}
Before the new one: say the Punjabi for to braid, then the Punjabi for a plait.
\end{quote}

\subsection*{You'll want to know: \pa{ਪਰਾਂਦਾ}}
\textbf{\pa{ਪਰਾਂਦਾ}} — \emph{parāndā} — \textquotedblleft{}a hair tassel; coloured yarn tied into a plait\textquotedblright{}.

\begin{quote}
\pa{ਭੈਣ} \pa{ਦੀ} \pa{ਗੁੱਤ} \pa{ਵਿੱਚ} \pa{ਲਾਲ} \pa{ਪਰਾਂਦਾ} \pa{ਹੈ।}
\end{quote}

\begin{quote}
\emph{bhaiṇ dī gutt vich lāl parāndā hai.}
\end{quote}

\begin{quote}
There is a red tassel in my sister's plait.
\end{quote}

A masculine noun.

\begin{cousinweb}
\textbf{\pa{ਪਰਾਂਦਾ}} (\emph{parāndā}) is inherited: CDIAL sets it under Sanskrit \emph{prānta}, \textquotedblleft{}end, edge\textquotedblright{}, and lists Punjabi \emph{parā̃dā}, \textquotedblleft{}particoloured yarn used by women for tying up their hair\textquotedblright{} (CDIAL 8940).
\end{cousinweb}

\begin{grammarlens}[title={What you've built}]
One more word for saying what you do.
\end{grammarlens}

\subsection*{Your turn}
\begin{itemize}
\raggedright
  \item \emph{Say it:} \emph{parāndā}
  \item \emph{Say it:} \emph{parāndā}, once more
  \item \emph{Say it:} \emph{bhaiṇ dī gutt vich lāl parāndā hai.} — There is a red tassel in my sister's plait.
  \item \emph{Recall:} say the Punjabi for to braid, then the Punjabi for a plait, then say \emph{parāndā} again
  \item \emph{Write it:} copy \textbf{\pa{ਪਰਾਂਦਾ}} once, checking the whole word
\end{itemize}

\begin{tcolorbox}[breakable,colback=teal!4,colframe=teal!35!black,title={Before you move on}]
What does \pa{ਪਰਾਂਦਾ} mean? (\textquotedblleft{}A hair tassel; coloured yarn tied into a plait\textquotedblright{}.) Say it once more.
\end{tcolorbox}
\section[\texorpdfstring{\pa{ਲਟ} (laṭ)}{ਲਟ (laṭ)}]{\pa{ਲਟ} (laṭ) — a lock of hair}

\label{lesson:PA-C231-lat}



\begin{quote}
Before the new one: say the Punjabi for a plait, then the Punjabi for a hair tassel.
\end{quote}

\subsection*{You'll want to know: \pa{ਲਟ}}
\textbf{\pa{ਲਟ}} — \emph{laṭ} — \textquotedblleft{}a lock of hair\textquotedblright{}.

\begin{quote}
\pa{ਉਸ} \pa{ਦੀ} \pa{ਇੱਕ} \pa{ਲਟ} \pa{ਮੱਥੇ} \pa{ਉੱਤੇ} \pa{ਹੈ।}
\end{quote}

\begin{quote}
\emph{us dī ikk laṭ matthe utte hai.}
\end{quote}

\begin{quote}
A lock of her hair is on her forehead.
\end{quote}

A feminine noun.

\begin{cousinweb}
\textbf{\pa{ਲਟ}} (\emph{laṭ}) is inherited from a reconstructed \emph{*laṭṭa}, \textquotedblleft{}hanging, tangled\textquotedblright{}; CDIAL lists Punjabi \emph{laṭ}, \textquotedblleft{}lock of hair\textquotedblright{} (CDIAL 10918).
\end{cousinweb}

\begin{grammarlens}[title={What you've built}]
One more word for saying what you do.
\end{grammarlens}

\subsection*{Your turn}
\begin{itemize}
\raggedright
  \item \emph{Say it:} \emph{laṭ}
  \item \emph{Say it:} \emph{laṭ}, once more
  \item \emph{Say it:} \emph{us dī ikk laṭ matthe utte hai.} — A lock of her hair is on her forehead.
  \item \emph{Recall:} say the Punjabi for a plait, then the Punjabi for a hair tassel, then say \emph{laṭ} again
  \item \emph{Write it:} copy \textbf{\pa{ਲਟ}} once, checking the whole word
\end{itemize}

\begin{tcolorbox}[breakable,colback=teal!4,colframe=teal!35!black,title={Before you move on}]
What does \pa{ਲਟ} mean? (\textquotedblleft{}A lock of hair\textquotedblright{}.) Say it once more.
\end{tcolorbox}
\section[(dialogue)]{Looking back: from Waking Up Early to Combing and Plaiting}

\label{lesson:PA-R231-review-225-231}



\begin{quote}
Say the Punjabi for to stretch out and to wake (someone) up.
\end{quote}

\subsection*{Your turn}
Say these aloud:

\begin{itemize}
\raggedright
  \item to stretch out, to wake (someone) up, not busy, to lie down, in a hurry
  \item ready, a toothbrush, to shiver, to shut a door, a walk
  \item hard to do, to rub, lather, grime, a tap
  \item to wring, to dry something, to hang up, a rope, to wipe
  \item a knot, to tie, to tighten, a cord, a knotted cloth button
  \item roomy, to take off, to wrap, loose trousers, a hanging earring
  \item oil, to braid, a plait, a hair tassel, a lock of hair
\end{itemize}

\begin{tcolorbox}[breakable,colback=teal!4,colframe=teal!35!black,title={Before you move on}]
To stretch out? (\textbf{\pa{ਪਸਾਰਨਾ}}, \emph{pasārnā}.) To wake (someone) up? (\textbf{\pa{ਜਗਾਉਣਾ}}, \emph{jagāuṇā}.) Not busy? (\textbf{\pa{ਵਿਹਲਾ}}, \emph{vihlā}.) To lie down? (\textbf{\pa{ਲੇਟਣਾ}}, \emph{leṭṇā}.) In a hurry? (\textbf{\pa{ਕਾਹਲਾ}}, \emph{kāhlā}.) Ready? (\textbf{\pa{ਤਿਆਰ}}, \emph{tiār}.) A toothbrush? (\textbf{\pa{ਬੁਰਸ਼}}, \emph{burash}.) To shiver? (\textbf{\pa{ਕੰਬਣਾ}}, \emph{kambṇā}.) To shut a door? (\textbf{\pa{ਭੇੜਨਾ}}, \emph{bheṛnā}.) A walk? (\textbf{\pa{ਸੈਰ}}, \emph{sair}.) Hard to do? (\textbf{\pa{ਔਖਾ}}, \emph{aukhā}.) To rub? (\textbf{\pa{ਮਲਣਾ}}, \emph{malṇā}.) Lather? (\textbf{\pa{ਝੱਗ}}, \emph{jhagg}.) Grime? (\textbf{\pa{ਮੈਲ}}, \emph{mail}.) A tap? (\textbf{\pa{ਟੂਟੀ}}, \emph{ṭūṭī}.) To wring? (\textbf{\pa{ਨਿਚੋੜਨਾ}}, \emph{nichoṛnā}.) To dry something? (\textbf{\pa{ਸੁਕਾਉਣਾ}}, \emph{sukāuṇā}.) To hang up? (\textbf{\pa{ਟੰਗਣਾ}}, \emph{ṭaṅgṇā}.) A rope? (\textbf{\pa{ਰੱਸੀ}}, \emph{rassī}.) To wipe? (\textbf{\pa{ਪੂੰਝਣਾ}}, \emph{pūnjhṇā}.) A knot? (\textbf{\pa{ਗੰਢ}}, \emph{gaṇḍh}.) To tie? (\textbf{\pa{ਬੰਨ੍ਹਣਾ}}, \emph{bannhṇā}.) To tighten? (\textbf{\pa{ਕੱਸਣਾ}}, \emph{kassṇā}.) A cord? (\textbf{\pa{ਡੋਰੀ}}, \emph{ḍorī}.) A knotted cloth button? (\textbf{\pa{ਘੁੰਡੀ}}, \emph{ghuṇḍī}.) Roomy? (\textbf{\pa{ਮੋਕਲਾ}}, \emph{moklā}.) To take off? (\textbf{\pa{ਲਾਹੁਣਾ}}, \emph{lāhuṇā}.) To wrap? (\textbf{\pa{ਲਪੇਟਣਾ}}, \emph{lapeṭṇā}.) Loose trousers? (\textbf{\pa{ਸੁੱਥਣ}}, \emph{sutthaṇ}.) A hanging earring? (\textbf{\pa{ਝੁਮਕਾ}}, \emph{jhumkā}.) Oil? (\textbf{\pa{ਤੇਲ}}, \emph{tel}.) To braid? (\textbf{\pa{ਗੁੰਦਣਾ}}, \emph{gundṇā}.) A plait? (\textbf{\pa{ਗੁੱਤ}}, \emph{gutt}.) A hair tassel? (\textbf{\pa{ਪਰਾਂਦਾ}}, \emph{parāndā}.) A lock of hair? (\textbf{\pa{ਲਟ}}, \emph{laṭ}.)
\end{tcolorbox}
